\par\noindent\textbf{(1)} 将二体引力 \(-\mu\boldsymbol r/r^3\) 在参考圆轨道一阶Taylor展开，叠加共转系Coriolis与离心惯性项，得到 \(\ddot x-2n\dot y-3n^2x=0\)，\(\ddot y+2n\dot x=0\)。其有效潮汐矩阵只在径向出现 \(-3n^2x\)。
\par\noindent\textbf{(2)} 记初值 \((x_0,y_0,v_{x0},v_{y0})\)，解为 \(x(t)=(4-3\cos nt)x_0+\sin(nt)v_{x0}/n+2(1-\cos nt)v_{y0}/n\)，\(y(t)=y_0+6(\sin nt-nt)x_0-2(1-\cos nt)v_{x0}/n+(4\sin nt-3nt)v_{y0}/n\)。可代回两式与四个初值核对。
\par\noindent\textbf{(3)} 切向解中线性漂移项是 \(-6nt\,x_0-3t\,v_{y0}\)，要使其消失需 \(v_{y0}=-2nx_0\)。代回径向解得 \(x(t)=x_0\cos nt+(v_{x0}/n)\sin nt\)，轨道只具有周期相对振荡。
\par\noindent\textbf{(4)} 邻近Kepler圆轨道的平均角速度 \(n(r)=\sqrt{\mu/r^3}\)，求导 \(dn/dr=-3n/(2r)\)。半长轴微差导致长期相位差随 \(t\) 线性增长，正是CW切向漂移模；相对距离不再满足小量条件时需回到非线性二体轨道。
